\documentclass[11pt,a4paper]{article}
\usepackage[margin=2.4cm]{geometry}
\usepackage{amsmath,amssymb}

\newcommand{\E}{\mathbb{E}}
% ---------------------------------------------------------------------------
% Inserting the real figures: in the executed notebook, either right-click and
% save the rendered plot, or add  fig.savefig('figN.pdf', bbox_inches='tight')
% after the corresponding plt.subplots(...) block and re-run that cell.  Then
% add \usepackage{graphicx} above and replace the \plotph{caption}{source}
% block with:
%   \begin{figure}[htbp]\centering
%   \includegraphics[width=0.92\linewidth]{figN}\caption{caption}\end{figure}
% Plot ordinals below count rendered figures from the top of each notebook and
% are stable across re-execution; the section heading + printed plot title are
% the robust identifiers.
% ---------------------------------------------------------------------------
\newcommand{\plotph}[2]{%
\begin{figure}[htbp]\centering
\framebox[0.92\linewidth]{\parbox[c][10.5em][c]{0.86\linewidth}{\centering
\textbf{[Plot placeholder]}\\[0.4em] #2}}
\caption{#1}
\end{figure}}

\setlength{\parskip}{0.45em}
\setlength{\parindent}{0pt}

\title{\vspace{-1.2em}\Large\textbf{The Equation Behind the Quote}\\[0.3em]
\large Stochastic optimal control for options market making, and how
physics-informed neural networks make it solvable under realistic desk
conditions\vspace{-0.4em}}
\author{\normalsize Prototype summary for the structured equity derivatives desk}
\date{\normalsize \today}

\begin{document}
\maketitle
\vspace{-1.6em}

\section*{Executive summary}

Whenever the desk answers a request for quote, it is implicitly solving an
optimisation problem: a wider skew earns a better margin but fills less
often, a fill moves an inventory that carries risk against limits, and the
right answer therefore depends not only on the instrument being quoted but
on the entire state of the book and of the market at that moment. Desks
handle this today through an accumulation of sensible local rules ---
vol marks, inventory add-ons, per-line overrides --- which, precisely
because each was designed in isolation, cannot be consistent with one
another across lines, inventory states, and market conditions. The central
message of this note is that when one writes the quoting trade-off down
carefully as a problem in stochastic optimal control, the mathematics
hands back a single equation --- a partial differential equation of the
same family as the Black--Scholes equation the desk already lives with ---
whose solution generates \emph{every} quote, on every line, on both sides,
at every inventory level, from one internally consistent object.

The reason this observation has remained largely academic is not that the
equation is exotic but that, for a realistic book, it is
high-dimensional: it carries one coordinate per position line on top of
time, spot, and volatility, and the classical numerical methods that
solve such equations by laying down a grid see their cost multiply with
every added line until the computation becomes impossible. The second
message of this note is that this particular obstacle has recently become
tractable. A class of solvers known as \emph{physics-informed neural
networks} --- explained from first principles in
Section~\ref{sec:pinn}, since the name is more intimidating than the
idea --- replaces the grid with a trained function, and in doing so
sidesteps exactly the cost explosion that ruled the approach out. We have
built a working prototype that solves the quoting equation for a small
option book under genuinely desk-like conditions --- per-line position
limits, and Greeks that move with spot and can flip sign, as they do
around barriers and autocall triggers --- and we propose a contained
pilot in which the same machinery is calibrated to the desk's own RFQ
history and run in shadow mode alongside the incumbent quoting logic.

\section{The quoting decision, as the desk lives it}

Consider what actually happens when a client request arrives for some
line $i$ in size $z_i$. The trader must choose a skew $\delta$ around the
model mid, knowing three things at once. First, the client's willingness
to trade falls away as the quote worsens, so that the probability of a
fill is some decreasing curve $\Lambda_i(\delta)$ of the skew --- a curve
which, importantly, is not a matter of theory but of measurement, since
the desk's own hit-and-miss history contains it. Second, a fill changes
the inventory, and inventory is not free: it consumes limit capacity, it
exposes the book to the market factors the desk cannot or does not hedge,
and, when the desk's real-world view of volatility differs from the
pricing measure, it also carries an expected drift, favourable or not.
Third, and most subtly, the value of taking on a given risk depends on
what the book already holds, which means that the correct skew on line
$i$ depends on the positions in every other line --- a dependence that
per-line rules, by construction, cannot express.

The control formulation takes these three ingredients exactly as stated
and asks for the quoting policy that maximises the risk-adjusted expected
profit over a trading horizon,
\begin{equation}
\sup_{\delta}\;\E\!\int_0^T\!\Big[\;
\underbrace{\textstyle\sum_{i,\pm} z_i\,\delta_i^{\pm}\,
\Lambda_i(\delta_i^{\pm})}_{\text{expected spread capture}}
\;+\;\underbrace{c(t,S,\nu)^{\!\top} q}_{\text{inventory carry}}
\;-\;\underbrace{\tfrac{\gamma}{2}\,
q^{\top}\Sigma(t,S,\nu)\,q}_{\text{inventory risk penalty}}
\;\Big]\,dt ,
\label{eq:obj}
\end{equation}
where $q$ is the vector of positions, one coordinate per line, the
penalty matrix $\Sigma$ encodes the residual risk of holding them, and
the supremum runs over all admissible quoting policies, with the desk's
position limits defining what admissible means. Nothing in this
expression is a modelling flourish; each term is a quantity the desk
already reasons about, merely written in one place.

\section{Where the partial differential equation appears}
\label{sec:pde}

The step from the objective~\eqref{eq:obj} to an equation is the dynamic
programming principle, and it is worth walking through slowly because it
is where the entire method acquires its structure. Define the value
function $v(t,S,\nu,q)$ to be the best risk-adjusted profit still
achievable between the present moment and the horizon, given the current
time, the current spot and volatility, and the current inventory vector.
Over the next instant, three things can happen to this value. Time
passes and the market state diffuses, which contributes the familiar
second-order terms --- in fact precisely the same diffusion operator,
call it $\mathcal L_{S,\nu}$, that drives the Black--Scholes and Heston
pricing equations the desk uses every day. The book accrues its running
carry and pays its running risk penalty on the inventory it currently
holds. And, with some intensity, a request arrives on one of the lines,
the desk quotes optimally against that client's fill curve, and if the
quote fills the inventory jumps by one trade. Requiring that these three
effects balance --- which is all the dynamic programming principle says
--- produces the Hamilton--Jacobi--Bellman equation
\begin{equation}
0 \;=\; \frac{\partial v}{\partial t}
\;+\;\mathcal L_{S,\nu}\,v
\;+\;c^{\!\top}q\;-\;\tfrac{\gamma}{2}\,q^{\top}\Sigma\,q
\;+\;\sum_{i,\pm}\;
\mathbf 1_{\{\text{limits}\}}\,
\sup_{\delta}\;\Lambda_i(\delta)\Big[\,z_i\,\delta
-\big(v(t,S,\nu,q)-v(t,S,\nu,q\mp z_i e_i)\big)\Big],
\label{eq:hjb}
\end{equation}
with the terminal condition that $v$ vanishes at the horizon. This is
the equation behind the quote. It is a partial differential equation in
the honest sense --- the unknown is a \emph{function} of time, spot,
volatility, and every inventory coordinate simultaneously --- and its
diffusion part is machinery the desk already trusts, while the genuinely
new ingredients are the inventory coordinates and the final sum, in
which each term is nothing other than the quote optimisation itself,
embedded inside the equation.

The reward for solving it is considerable, because once $v$ is known the
optimal quote on any line, on either side, in any state, is read off
pointwise: writing $p_i$ for the value difference per contract,
\begin{equation}
p_i \;=\; \frac{v(t,S,\nu,q)-v(t,S,\nu,q\mp z_i e_i)}{z_i},
\qquad
\delta_i^{*}\;=\;\arg\max_{\delta}\ \Lambda_i(\delta)\,(\delta-p_i),
\label{eq:quote}
\end{equation}
every quote decomposes into an inventory-dependent \emph{internal
transfer price} $p_i$ --- what this marginal risk is worth to this
particular book in this particular state --- plus an optimal margin
charged against the client's measured fill curve. This decomposition is
what makes the approach auditable rather than a black box: a trader or a
model-risk reviewer can inspect the two pieces separately, and both trace
back to parameters with direct empirical meaning.

A concrete and well-studied instance of this formulation, which we adopt
as our worked example, is the single-underlying option market-making
model of Baldacci, Bergault and Gu\'eant~\cite{bbg}: one stock with
Heston stochastic volatility, a book of listed calls that is
delta-hedged continuously so that vega remains as the residual risk, and
logistic fill curves calibrated in implied-volatility space so that a
quote at mid fills roughly a third of the time while a quote one vol
point through mid fills about seventy percent of the time --- two
anchors that a desk can estimate directly from its logs. The
formulation even repays the effort with structural insights of its own;
to give one example that we verified in simulation, the entire choice of
delta-hedging style, from plain delta-neutral to the minimum-variance
smile-adjusted delta, enters the quoting problem through nothing more
than a single scalar multiplier on the vega penalty, so that a debate the
desk might conduct at length collapses, within this framework, to one
coefficient.

To make their stylised case solvable with classical numerics, the
authors of~\cite{bbg} introduce two simplifying assumptions: the Greeks
of every line are frozen at their inception values, and the desk's risk
limit is a single band on the \emph{aggregate} vega of the book rather
than a limit per line. Under precisely those two assumptions the
inventory vector matters only through one number, the whole equation
collapses to three dimensions, and a conventional grid solver handles it
comfortably. The purpose of this note begins where those assumptions
end.

\section{Why realistic conditions defeat classical solvers}

Neither simplification survives contact with a real desk, and it is
worth being precise about what breaks, because the two failures are of
different kinds.

The first assumption to fall is the aggregate risk band, because actual
mandates impose position limits line by line, and the moment the
admissible set is a box rather than a band, the value of the book
genuinely ceases to be a function of its net Greek alone: two inventories
with identical net vega but different gross compositions sit at
different distances from different limits, face different feasible
unwinding paths, and are therefore worth different amounts. There is no
smaller equation hiding inside the problem; the inventory vector must be
carried whole, one coordinate per line.

The second assumption to fall is the frozen Greek, because a structured
book's vega moves with spot and volatility and, in the neighbourhood of
barriers and autocall triggers, changes sign --- so that the very
direction in which a new trade adds or reduces risk depends on where the
market is. In our prototype the minimal laboratory instance of this
behaviour is a tight call spread, long vega below its strikes and short
vega above, with the flip sitting inside a normal day's trading range;
any quoting logic keyed to a risk snapshot taken this morning will skew
in the wrong direction on the far side of that flip. Carrying live
Greeks means carrying spot itself as a state variable of the equation.

Put the two together and count coordinates: time, spot, volatility, and
one inventory dimension for every line the desk quotes. A classical
finite-difference solver must lay a grid over that entire space and
visit every node at every time step, so that its cost multiplies by the
grid width with each line added; a book of even seven or eight lines
already lies far beyond reach, and a structured book is out of the
question. This, and only this, is what kept the control formulation in
the academic literature: not any doubt about the equation, but the
absence of a way to solve it at realistic dimension.

\section{What a physics-informed neural network is}
\label{sec:pinn}

Since the term will be unfamiliar, it is worth building the idea up from
nothing, because the idea itself is simple. A neural network, stripped
of its mystique, is a flexible formula: a function $v_\theta(t,S,\nu,q)$
built from a few layers of elementary operations, whose many internal
parameters $\theta$ can be tuned until the formula behaves in whatever
way one requires. In the familiar applications of machine learning the
tuning target is data --- one shows the network examples of inputs and
correct outputs and adjusts the parameters until it reproduces them. A
\emph{physics-informed} neural network uses the same flexible formula
and the same tuning machinery, but replaces the data with an
\emph{equation}: since we know that the true value function satisfies
the HJB equation~\eqref{eq:hjb} at every point of the state space, we
can take the candidate formula $v_\theta$, substitute it into the
left-hand side of the equation at any state we please, and read off a
number --- the \emph{residual} --- that measures how badly the candidate
fails to be a solution at that point. Training then consists of sampling
states at random --- a time, a spot, a volatility, an inventory vector
--- evaluating the residual at each, and nudging the parameters to drive
the average squared residual toward zero, over and over, until the
formula satisfies the equation essentially everywhere it is asked. The
derivatives that the equation requires are not approximated by finite
differences but computed exactly from the formula itself by automatic
differentiation, the same mechanism that produces Greeks from modern
pricing libraries; and the terminal condition is built into the
architecture of the formula, so it holds by construction rather than by
training.

The decisive property of this procedure, for our purposes, is what it
does \emph{not} require: at no point is a grid ever constructed. The
network is trained by sampling points, not by enumerating them, and a
point in twelve dimensions is exactly as cheap to sample as a point in
three; the cost of the method therefore grows with the genuine
complexity of the value function being represented rather than with the
volume of the space it lives in, which is precisely the exchange one
needs in order to escape the dimensional explosion of
the previous section. Two further properties matter operationally. Once
trained, the network is just a formula, so evaluating a full set of
quotes for the entire book at the current state --- via
equation~\eqref{eq:quote} --- takes a fraction of a millisecond, which
is compatible with quoting latency in a way that no on-the-fly
computation could be. And because the object produced is the value
function itself rather than a table of quotes, everything the
formulation promises --- transfer prices, limit-aware fading, unwinding
urgency --- is available from the same trained object without further
work.

\section{The prototype: solving the realistic case}

We applied this solver to the quoting equation with \emph{both}
simplifying assumptions removed, on a deliberately small but
structurally honest book: three vanilla calls together with the tight
call spread described above, per-line position limits defining the
admissible inventory box, live Greeks supplied to the equation from a
gridded surrogate over spot and volatility fed by the same Heston pricer
throughout, and a trading horizon long enough that spot genuinely
traverses the spread's sign flip. The state of the equation is
therefore time, spot, volatility, and four inventory coordinates; the
network was given the raw state and nothing else --- in particular, no
aggregate-risk feature and no hint of the book's geometry --- and was
trained purely on the residual of equation~\eqref{eq:hjb}.

The trained object behaves the way the desk would want it to. Its quotes
respond to inventory in the correct direction on every line and fade
coherently as positions approach their limits, since the transfer price
in~\eqref{eq:quote} inherits the limit geometry directly from the
equation. Most tellingly, at an inventory long of the call spread, the
model's skew on that line \emph{reverses} as spot crosses the vega flip:
below the flip the book's spread position adds vega and the model quotes
to shed it, above the flip the same position reduces vega and the
quoting pressure changes sign --- exactly the behaviour that no
snapshot-based rule can produce, delivered here as an automatic
consequence of carrying spot inside the equation. As to whether the
answers can be trusted, we verified the solver in two complementary
ways, which we state briefly: when the two textbook simplifications are
quietly restored, the network reproduces the classically computable
solution --- at the level of the quotes themselves, to hundredths of a
cent --- and shows no spurious dependence on spot despite having spot as
an input; and in the realistic configuration, where no classical
benchmark exists at any book size, a martingale-based consistency check
--- simulating the book under the model's own policy and comparing the
realised risk-adjusted profit with the model's valuation --- closes to
within one standard error.

\plotph{Model bid and ask skews for a representative line across the
inventory range, showing the consistent inventory response and the
refusal to quote the breaching side at the position limit.}
{\emph{Grab from:} \texttt{stage2\_stage3\_ladder\_and\_frontier.ipynb},
Section~2, the code cell directly under the heading ``\emph{Quote-level
check vs the exact lattice}'' --- single-panel figure titled
``stage-2 quotes vs exact lattice'' (3rd plot of the notebook; x-axis
$n_1$, lattice lines + PINN markers; quotes end where the box limit
blocks the side).\\[0.3em]
\emph{Alternative (aggregate-band variant):}
\texttt{stage1\_pinn\_train\_and\_analysis.ipynb}, Section~4d, the
two-panel figure with titles ``K=10, T=1.0''\,/\,``K=12, T=1.0''
(5th plot of that notebook).}

\plotph{Left: the call spread's vega across spot, changing sign inside
the trading band. Right: the model's quote skew on the \emph{vanilla}
lines at a long-spread inventory, reversing as spot crosses the flip ---
holding the flip instrument, the desk sheds it on both sides; what
reverses is the direction of delta being shed, and therefore the quoting
of everything else on the book.}
{\emph{Grab from:} \texttt{stage2\_stage3\_ladder\_and\_frontier.ipynb},
Section~3. \emph{Left panel:} the first stage-3 code cell's figure,
titled ``live Greeks: the spread vega flips sign in the traded band''
(4th plot of the notebook; for the handout, keeping only the bold spread
line is cleanest). \emph{Right panel:} the notebook's 3b cell currently plots the
\emph{spread's own} skew, which does not reverse (you shed what you
hold on both sides of the flip); change one line in that cell --- plot
channels $1$ and $5$ (the ATM call) instead of $3$ and $7$ at the
long-spread inventory --- to produce the correct reversing quantity.
The exact-lattice verification of this corrected statement, with the
crossing at the flip in both BS and local-vol models, ships in the
companion \texttt{spotmm} repository (notebook
\texttt{spotmm\_full\_problem}, \S3).}

\section{What the desk gains}

The practical case for the approach rests on four properties that follow
from its construction rather than from tuning. The first is consistency:
because every quote on every line is generated from one value function,
the skews, the limit behaviour, the unwinding pressure and the carry
monetisation cannot contradict one another, in exactly the way that a
stack of independently maintained rules eventually must. The second is
auditability: the decomposition~\eqref{eq:quote} of each quote into a
transfer price and a margin gives trading and model risk two separately
inspectable quantities, and every parameter feeding them --- the fill
curves, the request intensities, the risk penalty --- has a direct
empirical counterpart, with the fill curves in particular calibratable
from the desk's own RFQ hit-and-miss history rather than from any
theoretical assumption. The third is extensibility: the prototype's
live-Greek layer is simply a surrogate over spot and volatility fed by a
pricer, and the identical pattern accepts the desk's own structured
pricers, so that the call spread of the prototype is the minimal
stand-in for the barrier and autocall vega profiles the desk actually
runs. The fourth is deployability: the intended mode of operation is
decision support --- model quotes and transfer prices surfaced alongside
the trader's own judgement, with post-fill markout monitoring for
adverse selection --- rather than autonomous execution, so that the
system earns trust in the same way a junior trader would.

\section{Scope, limitations, and a proposed pilot}

In the interest of a fair pitch, the current prototype should be
described exactly as it is. It runs on a single underlying with
unit-size requests; its neural trainings were performed at deliberately
modest computational budgets, with the identical code scaling to full
budgets on a single data-centre GPU; the hedge is assumed frictionless
and continuous, and introducing hedge costs or caps is a genuine
extension rather than a switch; and the Greeks, while fully spot- and
vol-dependent, are held fixed in time over the short demonstration
horizon. None of these is a limitation of the formulation --- each is a
boundary of what has so far been built and tested, drawn deliberately so
that every claim above corresponds to something that was actually run.

The natural next step is a contained pilot on the desk's own terms.
First, calibrate the per-line fill curves and request intensities from
the desk's RFQ logs, in implied-volatility space and by client tier if
the data supports it, since these curves are the empirical heart of the
model. Second, configure a pilot sub-book --- for instance a set of
index vanillas together with one structured overlay line --- with the
desk's actual limits as the admissible set, and train the solver at full
budget on that configuration. Third, run in shadow: surface the model's
quotes and transfer prices alongside the incumbent logic, and monitor
markouts and limit utilisation. Fourth, evaluate head to head under
common random scenarios, comparing realised spread capture against
inventory-risk variance, the coherence of skews across lines, and
behaviour as limits are approached. Each step produces something of
standalone value --- the calibrated fill curves alone are informative
about the franchise --- and no step commits the desk beyond the one
before it.

\plotph{Pilot evaluation mock-up: distribution of risk-adjusted P\&L,
model versus incumbent skew logic, under common random scenarios.}
{\emph{To be produced during the pilot} --- no notebook source exists
yet. \emph{Stand-in, if one is wanted for the handout:} the repository
simulator already produces this style of comparison; run
\texttt{sim/lockstep.simulate\_atoms} under two quote-table sets on
common random numbers (the optimal vs.\ perturbed-policy construction of
\texttt{tests/test\_corrector.py}) and histogram the two per-path
risk-adjusted reward vectors side by side.}

\begin{thebibliography}{9}\small
\bibitem{as} M.~Avellaneda, S.~Stoikov,
\emph{High-frequency trading in a limit order book},
Quantitative Finance 8(3), 2008.
\bibitem{gueant} O.~Gu\'eant,
\emph{The Financial Mathematics of Market Liquidity: From Optimal
Execution to Market Making}, Chapman \& Hall/CRC, 2016.
\bibitem{bbg} B.~Baldacci, P.~Bergault, O.~Gu\'eant,
\emph{Algorithmic market making for options}, arXiv:1907.12433.
\bibitem{pinn} M.~Raissi, P.~Perdikaris, G.E.~Karniadakis,
\emph{Physics-informed neural networks: A deep learning framework for
solving forward and inverse problems involving nonlinear partial
differential equations}, Journal of Computational Physics 378, 2019.
\end{thebibliography}

\end{document}
